% --- 仅输入正文内容 ---

\section{实验目的：}
\begin{enumerate}[label=\arabic* 、, leftmargin=4em, labelsep=0pt]
    \item 学习可见分光光度计的基本原理：
    \item 掌握液体吸光度测量方法，利用吸光度与液体浓度的关系测量液体浓度：
    \item 掌握测量不同颜色物体的反射和透射光谱，并辨别颜色。
\end{enumerate}

\section{实验原理：}

\subsection{分子能级与光谱}
分子一般包含若干个原子核和电子，其能量可以表示为：$E = E_e + E_v + E_r + E_n + E_t$。其中，$E_e$、$E_v$、$E_r$ 分别表示电子、原子间振动、分子转动能量。分子产生跃迁所吸收能量的辐射频率直接依赖于能级的差距。分子的紫外可见吸收光谱是由电子跃迁引起的，红外吸收光谱是由于分子中基团的振动和转动能级跃迁引起的。

% (此处可插入图片环境)
% \begin{figure}[htbp]
%   \centering
%   \includegraphics[width=0.6\textwidth]{path_to_image}
%   \caption{分子典型能级图}
% \end{figure}

\subsection{光谱学概念}
\begin{enumerate}[label=\arabic* 、, leftmargin=4em, labelsep=0pt]
    \item \textbf{反射率}：指的是反射波与入射波的功率比值，即 $R(\lambda) = I_r / I_o$；
    \item \textbf{透过率}：指的是透射波与入射波的光强比值，即 $T(\lambda) = I_t / I_o$；
    \item \textbf{吸光度}：衡量物质对光的吸收程度。
    \item \textbf{色品图}：表示颜色系统的工具，可直观了解和比较不同颜色之间的关系和特性。
\end{enumerate}

\subsection{朗伯比尔定律}
指当一束平行单色光通过均匀、非散射的稀溶液时，溶液对光的吸收程度与溶液的浓度及液层厚度的乘积成正比。

朗伯比尔定律描述了浓度与吸光度的关系：
$A = K C L$
其中 K 为波长处的样品消光系数，C 为样品的浓度，L 为吸收光程。当光程一定时，吸光度A与浓度c呈线性关系。同时，与入射光和出射光强度有关：$A = -\ln(T)$。

\section{实验仪器：}
\subsection{试剂与材料}
酒精、内酮、不同浓度日落黄染料、不同浓度苋菜红染料、不同颜色色卡等。

\subsection{仪器}
光源、光谱仪、光纤、积分球、标准白板、比色皿、计算机等。
\begin{enumerate}[label=\arabic* 、, leftmargin=4em, labelsep=0pt]
    \item \textbf{光源}：光源的作用是提供激发能，发出足够强的连续光谱；
    \item \textbf{光谱仪}：把光源发出的光分离出所需要的单色光。通常由入射狭缝、准直镜、色散元件（光栅）、物镜和出口狭缝构成。
    \item \textbf{比色皿}：用于盛放试液。石英比色皿用于可见光区的测量，玻璃池只用于可见光区。
\end{enumerate}

\section{实验内容：}
\subsection{液体吸光度}
\begin{enumerate}[label=\arabic* 、, leftmargin=4em, labelsep=0pt]
    \item 将纯水倒入比色皿中作为对照液体，将比色皿放入比色皿支架；
    \item 使用一根光纤连接比色皿支架与光谱仪，另一根光纤连接比色皿支架与光源；
    \item 通过 USB 数据线连接光谱仪至电脑，用电源线连接光源和插座；
    \item 双击电脑桌面上 “Oceanview” 软件，配置吸收率扫描并调整积分时间校准等设置并保存黑暗参考。
    \item 放置样品进行测量读取曲线。
\end{enumerate}

\subsection{测量物体颜色}
\begin{enumerate}[label=\arabic* 、, leftmargin=4em, labelsep=0pt]
    \item 拧开激光器上的螺帽，连接光纤与设备；
    \item 调整软件并使用积分球对准标准白板进行基础光学设置归一化。
\end{enumerate}

\section{注意事项：}
\begin{enumerate}[label=\arabic* 、, leftmargin=4em, labelsep=0pt]
    \item 在吸光度实验中，调整软件积分时间前应先将激光器强度调至最大，再调整积分时间，使激光强度处于合适范围。
    \item 在吸光度实验中，积分时间与多次扫描平均值的乘积应处在合适的范围内。过小会导致测量不准确，过大会导致测量过慢，影响实验。
    \item 积分球下端黄色薄膜需保存好，做完实验需重新贴上，以免积分球损坏。
    \item 操作激光时应注意安全，请勿将光纤口对准眼睛。
\end{enumerate}

\section{实验数据：}
% 实验数据表格
\begin{table}[htbp]
    \centering
    \caption{颜色体系记录与浓度值}
    \begin{tabular}{llcc}
        \hline
        颜色类型 & 样品浓度 & 特征峰位 (nm) \\
        \hline
        红色 & 1.0 $\times 10^{-4}$ & 655 \\
        黄色 & 1.0 $\times 10^{-4}$ & 550 \\
        绿色 & 1.0 $\times 10^{-4}$ & 515 \\
        浅蓝色 & 1.0 $\times 10^{-4}$ & 480 \\
        深蓝色 & - & 413 \\
        紫色 & - & 425 \\
        \hline
    \end{tabular}
\end{table}

\section{数据分析：}
从图 7.1 可知，在可见光范围内，白光的反射率为 100\%，不同颜色需要在指定波长范围内才可 100\% 反射，因此能够分析出不同色卡其本身的颜色反射出的波长，并以此计算其在色品图中的具体位置：红色：655nm；黄色：550nm；绿色：515nm；浅蓝色：480nm；深蓝色：413nm；紫色：425nm。

从图 7.2 可知，不同物质（溶液）具有不同的吸光度和特征峰，可以通过测定某种物质的吸收谱，根据峰值和峰位来判断出物质或种类。例如实验中用到的苋菜红溶液的特征吸收峰位为 465nm，对该波长的光的吸收度为 0.23；日落黄溶液的特征吸收峰位为 490nm，对该波长的光的吸收度为 0.26。

从图 7.3 可知，随着日落黄溶液的浓度的升高，相应溶液的吸光度增加，可以认为是因为溶液中溶质的增加，光的透过率下降，即吸光度增加。原则上同一物质的特征吸收峰的峰位是不会随着浓度的变化而变化，但由于浓度之间的数量级差异，出现吸收饱和，因此从图像中可以看出来峰位随浓度的升高向长波方向偏移。从图像上反馈为峰下侧面积增大。测试样品的浓度不宜过大，否则偏移现象会导致定性错误。

从图 7.4 可知，将吸光度进行对数化，与浓度进行拟合，可知在较低浓度时，样品的吸光度随着浓度的升高而升高，这符合实验事实规律，即朗伯比尔定律。但同样由于样本浓度太高，测试时接近吸收饱和，导致曲线并没有完全呈线性变化。

\section{思考题：}
\begin{enumerate}[label=\arabic* 、, leftmargin=4em, labelsep=0pt]
    \item \textbf{紫外可见光谱的范围一般分为三个区域：远紫外区、近紫外区和可见光区那么这三个区域对应的波长范围分别是多少？} \\
    答；远紫外波长：10-200nm、近紫外波长：200-400nm、可见光：400-780nm。
    
    \item \textbf{在可见光区，颜色与波长有一定的对应关系。如果某溶液能够强烈地吸收480~510nm波长的光，那么该溶液会呈现什么颜色？如果某溶液呈橙红色，那么该溶液的可见光光谱在哪个波段会缺失？} \\
    答：1、溶液呈现出其互补色，为红色或紫色；2、说明吸收红色和橙色光谱，相对应的缺失其互补色为蓝色、绿色及两者混合色，对应光谱上425nm-570nm。
    
    \item \textbf{如果入射光为非单色光或不平行入射，吸收定律（郎伯比尔定律）会发生怎样的偏离？} \\
    答：1、非单色光：非单色光包含多种波长的光，不同波长的光能被不同程度地吸收。如果溶液对不同波长的光吸收不同，非单色光源中的某些波长可能会被强烈吸收而其他波长则不会，导致测量结果偏离。\\
    2、不平行入射：光在通过溶液时的路径长度可能会变化，导致不同部分的光受不同程度的吸收。不平行的光束可能导致溶液中某些区域接受到的光强比其他区域更强，影响吸光度的测量。
\end{enumerate}